% T-13. Plan de pruebas y validación.
% Índice propio del PDF separado; no se repite en un ensamblado integrado.
\ifdefined\FormularioActual
\makeatletter
\audit@iniciarpartes
\makeatother
{\renewcommand\textoEncabezado{\color{audit-marino}Índice}\indiceDetallado}
\clearpage
\makeatletter
\audit@marcasubdoc{Subdocumento 9}
\makeatother
\markboth{Formulario T-13}{}
\registrarFolio{formsd-9}{Formulario T-13}
\phantomsection
\addcontentsline{toc}{section}{Formulario T-13}
\fi
\begin{formulario}{T-13}
\begin{contenidoExigido}{T-13}
\exigencia{Niveles y tipos de prueba conforme a ISO/IEC/IEEE 29119.}{\hyperref[sec:9-t13-1]{Apartado 1}, folio \pageref{sec:9-t13-1}}
\exigencia{Ambientes y datos de prueba.}{\hyperref[sec:9-t13-2]{Apartado 2}, folio \pageref{sec:9-t13-2}}
\exigencia{Criterios de entrada y de salida.}{\hyperref[sec:9-t13-3]{Apartado 3}, folio \pageref{sec:9-t13-3}}
\exigencia{Automatización de pruebas.}{\hyperref[sec:9-t13-4]{Apartado 4}, folio \pageref{sec:9-t13-4}}
\exigencia{Calendario de las pruebas de carga, resiliencia, recuperación ante desastres y seguridad ofensiva.}{\hyperref[sec:9-t13-5]{Apartado 5}, folio \pageref{sec:9-t13-5}}
\end{contenidoExigido}

\subsection*{Niveles y tipos de prueba conforme a ISO/IEC/IEEE 29119.}\label{sec:9-t13-1}
\addcontentsline{toc}{subsection}{Niveles y tipos de prueba conforme a ISO/IEC/IEEE 29119.}
Se aplica planificación, seguimiento, diseño, preparación, ejecución y cierre de los procesos de prueba \parencite{iso29119_2_2021}. La estructura de nueve campos del catálogo es una convención de audIT. Se diseñan 25 casos unitarios, 25 de integración, 20 de sistema, 12 de rendimiento y 8 de seguridad, 15 de hardware/HIL y 15 UAT: 120 identificadores únicos. Unitarias contrastan reglas/oráculos; integración valida esquemas e idempotencia; sistema verifica flujos completos; no funcionales miden carga, seguridad y recuperación; HIL utiliza dispositivos homologados; UAT comprueba tareas con contraparte. Véase \verSeccionFolio{sec:9-estrategia}.

\subsection*{Ambientes y datos de prueba.}\label{sec:9-t13-2}
\addcontentsline{toc}{subsection}{Ambientes y datos de prueba.}
CI utiliza dobles y fixtures sintéticos; QA integra contratos; staging reproduce configuración/carga; HIL conserva modelo, firmware, señales y permisos; UAT usa terminales y usuarios designados con autorización. Producción se utiliza solo en ejercicio autorizado, con respaldo y reversión. Se versionan entradas, semilla/reloj, oráculo y manifiesto. No se atribuyen marcas o API ensayadas al parque existente sin homologación. Datos operacionales requieren acceso autorizado y minimización. Véase \verSeccionFolio{sec:9-ambientes}.

La homologación física utiliza iWave G26I con audIT EdgeHub, lector CAN sin contacto Technoton CANCrocodile y las balizas Bluetooth de la innovación 3. Se distinguen los 182 equipos previstos de los 192 terceros con dispositivos existentes: sus plataformas requieren acceso autorizado y no se presume API ni exportación disponible. El catálogo T-17 verifica lectura, integridad, datos ausentes, antigüedad y conciliación, sin atribuir al parque real el resultado de dobles sintéticos. Para frío, rango y precisión provienen del modelo y la homologación; no se trasladan especificaciones de sondas PT100.

\subsection*{Criterios de entrada y de salida.}\label{sec:9-t13-3}
\addcontentsline{toc}{subsection}{Criterios de entrada y de salida.}
Entrada: versión/configuración y ambiente identificados, requisito/caso/criterio fijados, datos y permisos disponibles, observabilidad y reversión verificadas. Salida: todos los criterios aplicables cumplidos, reporte reproducible y defectos relacionados, sin críticos/altos para cierre de marcha blanca. Se aplican los mínimos de \verSeccionFolio{sec:9-minimos-trazados}; cuatro semanas sostenidas y las seis condiciones acumulativas de \verSeccionFolio{sec:9-condiciones-contractuales} gobiernan aceptación. Un promedio o cobertura de código no reemplaza otro criterio.

Los controles siguientes complementan los criterios contractuales de S9 §9.3.6. La \tabref{9-t13-gates} relaciona umbral, método y evidencia; los controles físicos son compromisos de homologación de la configuración ofertada, no certificaciones ya obtenidas. audIT provee o contrata el banco y los instrumentos, sin presumir infraestructura del mandante.


\begin{tabla}[ancho=completo,fuente=condensada]{Controles bloqueantes de software y homologación física}{9-t13-gates}{L{27mm} X X X}
Control & Umbral y alcance & Método y caso & Evidencia de salida \\ \midrule
Cobertura & Líneas de negocio \(\ge\)70 \% contractual; ramas \(\ge\)80 \% adicional & GitLab Runner y reporte compatible con SonarQube; catálogo unitario CP-UNIT-01–25 & Versión, numeradores, denominadores, exclusiones y cero pruebas fallidas \\
Complejidad & \(\le\)15 por función de negocio, compromiso audIT & SonarQube y revisión de descomposición; S9 §9.1.1 & Reporte estático y cierre de funciones fuera de límite \\
Vulnerabilidades & Cero críticas o altas abiertas en el artefacto promovido & SonarQube, Trivy y OWASP ZAP; CP-SEC-01–08 según alcance & Informes ligados a huella, reevaluación y SBOM \\
Lectura CAN pasiva & Pérdida <0,1 \%; cero tramas transmitidas por audIT & CANoe y analizador independiente; CP-HW-04 & Contadores, captura, autorización y mapa de señales \\
Energía en reposo & <50 mA tras 30 min, compromiso audIT & Instrumento de corriente sobre conjunto alimentado; CP-HW-05 & Serie de corriente, tensión, duración y configuración de periféricos \\
Ambiente térmico y vibración & Operación entre -20 °C y +70 °C; perfil SAE J1455 aprobado para montaje & Cámara y mesa de vibración; CP-HW-07 & Protocolo identificado, certificado de calibración e informe funcional \\
Protección exterior & IP67 del conjunto expuesto: gabinete, conectores y cableado exterior & Laboratorio y procedimiento IEC 60529; CP-HW-08 & Configuración de montaje, informe de polvo/agua y verificación funcional \\
Actualización de firmware & Solo detenido en terminal autorizado; recuperación A/B sin pérdida & Control de ubicación/estado y fallos de imagen/energía; CP-SYS-20 y CP-HW-09 & Rechazos fuera de ventana, huellas, arranque, reversión y conciliación \\
\end{tabla}

Cada informe conserva la versión del procedimiento, modelo, montaje y calibración aplicables. Una ficha de fabricante o un resultado sobre otro montaje no sustituye la homologación del conjunto. Un fallo impide liberar ese equipo o incremento hasta su corrección y reevaluación. El perfil térmico y vibratorio se documenta antes del ensayo; los números de un fixture no se atribuyen a una cláusula normativa sin respaldo.

\subsection*{Automatización de pruebas.}\label{sec:9-t13-4}
\addcontentsline{toc}{subsection}{Automatización de pruebas.}
Cada cambio ejecuta compilación, pruebas unitarias y regresión; contratos e integración se ejecutan contra dobles versionados y en QA. La cobertura automatizada de líneas ejecutables de lógica de negocio \(\ge\)70 \% es el mínimo contractual bloqueante (RT-04.11); la cobertura de ramas \(\ge\)80 \% es un compromiso adicional bloqueante de audIT, coherente con S6 Tabla 6.3. Se calculan líneas cubiertas/líneas ejecutables y ramas cubiertas/ramas instrumentadas por separado, con módulos y exclusiones justificadas; no se incluyen dependencias o código generado para elevar el resultado. Análisis estático/dependencias y pruebas de autorización acompañan cada promoción; el resultado se conserva por versión. Carga, cortes de enlace y recuperación requieren script, manifiesto y cronómetro; HIL/UAT conservan las verificaciones manuales que no puede suplir un pipeline. Un fallo bloquea promoción y dispara corrección/reevaluación; una excepción no rebaja mínimos contractuales ni los gates adicionales bloqueantes comprometidos. Véase sección 9.2.3.

\subsection*{Calendario de las pruebas de carga, resiliencia, recuperación ante desastres y seguridad ofensiva.}\label{sec:9-t13-5}
\addcontentsline{toc}{subsection}{Calendario de las pruebas de carga, resiliencia, recuperación ante desastres y seguridad ofensiva.}
El perfil de carga común es 380 sesiones nominales (80 internas, 150 conductores, 50 transportistas y 100 clientes) y 570 en estrés (120, 225, 75 y 150 respectivamente). Son hipótesis conservadoras de S4, no mediciones del Caso; CP-PERF-01 preserva esa mezcla y comprueba los umbrales contractuales a 1,5 veces el peak. Un cambio de perfil exige actualizar S9 y T-17 antes de ensayar.

Antes de ejecutar la carga se aprobará un manifiesto que identifique versiones de S2 y del perfil de arquitectura, poblaciones de referencia, proporciones de simultaneidad, operaciones por actor, tiempos de espera, rampa, duración, volumen de datos y ambiente. H1 contrastará las hipótesis con sesiones y actividad observadas; sin ese contraste, 380 y 570 siguen siendo escenarios sintéticos y el resultado no acredita representatividad de la operación real. El perfil incluye por separado la cota de viajes diarios de S2 y la concurrencia; no convierte una magnitud en la otra. Los umbrales contractuales se mantienen aunque la calibración obligue a ajustar la capacidad.

Construcción E1: A09/A10/A11; certificación E1 A12, M12/H5, con informes separados de las cuatro familias y UAT. Marcha blanca A15, M13--M15; producción A19, M16/H7. Construcción E2 A20 y certificación A22, M18/H10, incluyendo regresión E1 y las cuatro familias. Marcha blanca A24, M19--M20; aceptación final A25, M21/H12. DR y seguridad ofensiva requieren autorización previa y reversión; las ventanas contractuales no se desplazan por ensayos. Resiliencia y DR se repiten semestralmente en operación; intrusión anual y antes de producción por tercero independiente, con informe íntegro y remediación. Se realizan dos ensayos de migración antes de la definitiva (FEP02, 20.1, pp. 34--35; RT-11.20, p. 24). Véase \verSeccionFolio{sec:9-calendario-pruebas}.
\end{formulario}
